%%%
%%% Radical "⼸"
%%%
\section*{Radical 57: ``⼸''}\addcontentsline{toc}{section}{Radical 57: ⼸}\addcontentsline{loh}{figure}{\#\#\#\# 57: ⼸}

%%%%%%%%%% 弓 %%%%%%%%%%
\subsection*{弓}\addcontentsline{loh}{figure}{弓}

\begin{Entry}{弓}{3}{⼸}[Kangxi 57]
  \begin{Phonetics}{弓}{gong1}[][HSK 7-9]
    \definition*{s.}{Sobrenome: Gong}
    \definition{clas.}{uma antiga unidade de comprimento para medir a terra, igual a cinco 尺}
    \definition[张]{s.}{arco | qualquer coisa em forma de arco | Obsoleto: ferramenta de madeira de medição de terreno; divisores de madeira para medição de terrenos; arco de medição; régua escalonada (1,5m)}
    \definition{v.}{dobrar; arquear; curvar; entortar}
  \seealsoref{尺}{chi3}
  \end{Phonetics}
\end{Entry}

%%%%%%%%%% 引 %%%%%%%%%%
\subsection*{引}\addcontentsline{loh}{figure}{引}

\begin{Entry}{引}{4}{⼸}
  \begin{Phonetics}{引}{yin3}[][HSK 4]
    \definition*{s.}{Sobrenome: Yin}
    \definition{clas.}{uma unidade de comprimento equivalente a =33,333\dots metros}
    \definition{v.}{puxar; esticar | liderar; conduzir; guiar | sair; deixar | sobressair | atrair; provocar; trazer à existência | causar; provocar | citar; ser usado como evidência ou justificativa}
  \synonymref{避}{bi4}
  \synonymref{拉}{la1}
  \synonymref{离}{li2}
  \synonymref{牵}{qian1}
  \synonymref{去}{qu4}
  \synonymref{退}{tui4}
  \synonymref{拖}{tuo1}
  \synonymref{曳}{ye4}
  \synonymref{拽}{zhuai4}
  \antonymref{短}{duan3}
  \antonymref{进}{jin4}
  \antonymref{来}{lai2}
  \antonymref{留}{liu2}
  \antonymref{送}{song4}
  \antonymref{推}{tui1}
  \antonymref{抑}{yi4}
  \antonymref{诱}{you4}
  \antonymref{止}{zhi3}
  \end{Phonetics}
\end{Entry}

\begin{Entry}{引人入胜}{4,2,2,9}{⼸,⼈,⼊,⾁}
  \begin{Phonetics}{引人入胜}{yin3ren2-ru4sheng4}[][HSK 7-9]
    \definition{expr.}{conduzir alguém à parte interessante de algo; sedutor; atraente; cativante; emocionante; encantador; fascinante; que atraia muito alguém; interessante e absorvente (intrigante); levar alguém a uma nova perspectiva; abrir uma nova visão; deslumbrante; envolver o leitor em uma cena bela (referindo"-se à paisagem, à escrita, etc.)}
  \synonymref{扣人心弦}{kou4ren2xin1xian2}
  \end{Phonetics}
\end{Entry}

\begin{Entry}{引人注目}{4,2,8,5}{⼸,⼈,⽔,⽬}
  \begin{Phonetics}{引人注目}{yin3ren2-zhu4mu4}[][HSK 7-9]
    \definition{expr.}{impressionante; chamativo; que prende a atenção; significa descrever uma pessoa ou coisa que é muito peculiar e atrai a atenção das pessoas}
  \synonymref{举世瞩目}{ju3shi4-zhu3mu4}
  \synonymref{怡然自得}{yi2ran2-zi4de2}
  \antonymref{默默无闻}{mo4mo4-wu2wen2}
  \end{Phonetics}
\end{Entry}

\begin{Entry}{引入}{4,2}{⼸,⼊}
  \begin{Phonetics}{引入}{yin3ru4}[][HSK 7-9]
    \definition{v.}{introduzir; atrair; trazer de outro lugar | introduzir (pessoal, capital, tecnologia, equipamentos, etc.) de outras regiões ou países estrangeiros}
  \synonymref{引进}{yin3jin4}
  \synonymref{引用}{yin3yong4}
  \end{Phonetics}
\end{Entry}

\begin{Entry}{引发}{4,5}{⼸,⼜}
  \begin{Phonetics}{引发}{yin3fa1}[][HSK 7-9]
    \definition{v.}{despertar; desencadear; acionar; iniciar (ou acender, provocar)}
  \synonymref{激发}{ji1fa1}
  \synonymref{激励}{ji1li4}
  \end{Phonetics}
\end{Entry}

\begin{Entry}{引用}{4,5}{⼸,⽤}
  \begin{Phonetics}{引用}{yin3yong4}[][HSK 7-9]
    \definition{v.}{citar; transcrever; utilizar as palavras de outras pessoas ou eventos passados como evidência para apoiar e comprovar o próprio ponto de vista | nomear; recomendar; nomear alguém para um cargo ou para realizar um trabalho}
  \synonymref{引入}{yin3ru4}
  \end{Phonetics}
\end{Entry}

\begin{Entry}{引导}{4,6}{⼸,⼨}
  \begin{Phonetics}{引导}{yin3dao3}[][HSK 4]
    \definition{v.}{conduzir; guiar; liderar; andar na frente e deixar que os outros sigam atrás para ver ou andar; usar imagens ou sinais para mostrar às pessoas para onde ir | esclarecer; fornecer orientação em termos de ideias, métodos, conceitos, etc.}
  \synonymref{带领}{dai4ling3}
  \synonymref{启发}{qi3fa1}
  \synonymref{率领}{shuai4ling3}
  \synonymref{向导}{xiang4dao3}
  \synonymref{引领}{yin3ling3}
  \synonymref{指导}{zhi3dao3}
  \synonymref{指点}{zhi3dian3}
  \synonymref{指引}{zhi3yin3}
  \antonymref{误导}{wu4dao3}
  \antonymref{阻碍}{zu3'ai4}
  \antonymref{阻拦}{zu3lan2}
  \antonymref{阻挠}{zu3nao2}
  \antonymref{阻止}{zu3zhi3}
  \end{Phonetics}
\end{Entry}

\begin{Entry}{引进}{4,7}{⼸,⾡}
  \begin{Phonetics}{引进}{yin3jin4}[][HSK 4]
    \definition{v.}{importar; trazer de fora | recomendar; dar uma indicação}
  \seealsoref{引入}{yin3ru4}
  \synonymref{介绍}{jie4shao4}
  \synonymref{推荐}{tui1jian4}
  \synonymref{吸纳}{xi1na4}
  \synonymref{吸收}{xi1shou1}
  \synonymref{招揽}{zhao1lan3}
  \synonymref{招收}{zhao1shou1}
  \antonymref{驱逐}{qu1zhu2}
  \antonymref{输出}{shu1chu1}
  \antonymref{淘汰}{tao2tai4}
  \end{Phonetics}
\end{Entry}

\begin{Entry}{引经据典}{4,8,11,8}{⼸,⽷,⼿,⼋}
  \begin{Phonetics}{引经据典}{yin3jing1-ju4dian3}[][HSK 7-9]
    \definition{expr.}{``Citar capítulo e versículo.''; citar os clássicos; ciare obras de referência}
  \end{Phonetics}
\end{Entry}

\begin{Entry}{引诱}{4,9}{⼸,⾔}
  \begin{Phonetics}{引诱}{yin3you4}[][HSK 7-9]
    \definition{v.}{atrair; seduzir; tentar; induzir e orientar pessoas a fazer coisas ruins}
  \synonymref{勾结}{gou1jie2}
  \synonymref{迷惑}{mi2huo4}
  \synonymref{引导}{yin3dao3}
  \synonymref{诱惑}{you4huo4}
  \antonymref{导致}{dao3zhi4}
  \end{Phonetics}
\end{Entry}

\begin{Entry}{引起}{4,10}{⼸,⾛}
  \begin{Phonetics}{引起}{yin3qi3}[][HSK 4]
    \definition{v.}{causar; despertar; levar a; desencadear; dar origem a}
  \synonymref{导致}{dao3zhi4}
  \synonymref{激发}{ji1fa1}
  \synonymref{诱发}{you4fa1}
  \antonymref{避免}{bi4mian3}
  \antonymref{遏制}{e4zhi4}
  \antonymref{平息}{ping2xi1}
  \antonymref{消除}{xiao1chu2}
  \antonymref{抑制}{yi4zhi4}
  \end{Phonetics}
\end{Entry}

\begin{Entry}{引领}{4,11}{⼸,⾴}
  \begin{Phonetics}{引领}{yin3ling3}[][HSK 7-9]
    \definition{v.}{liderar; guiar | aguardar ansiosamente por algo; esperar ansiosamente}
  \end{Phonetics}
\end{Entry}

\begin{Entry}{引擎}{4,16}{⼸,⼿}
  \begin{Phonetics}{引擎}{yin3qing2}[][HSK 7-9]
    \definition[个,台]{s.}{motor, referindo"-se especificamente a motores térmicos como motores a vapor e motores de combustão interna | Empréstimo Linguístico: \emph{engine}}
  \end{Phonetics}
\end{Entry}

%%%%%%%%%% 弘 %%%%%%%%%%
\subsection*{弘}\addcontentsline{loh}{figure}{弘}

\begin{Entry}{弘}{5}{⼸}
  \begin{Phonetics}{弘}{hong2}
    \definition*{s.}{Sobrenome: Hong}
    \definition{adj.}{grande; grandioso; magnífico}
    \definition{v.}{ampliar; expandir | promover}
  \end{Phonetics}
\end{Entry}

\begin{Entry}{弘扬}{5,6}{⼸,⼿}
  \begin{Phonetics}{弘扬}{hong2yang2}[][HSK 7-9]
    \definition{v.}{melhorar; levar adiante; desenvolver e expandir; promover; promover vigorosamente}
  \end{Phonetics}
\end{Entry}

%%%%%%%%%% 弟 %%%%%%%%%%
\subsection*{弟}\addcontentsline{loh}{figure}{弟}

\begin{Entry}{弟}{7}{⼸}
  \begin{Phonetics}{弟}{di4}[][HSK 1]
    \definition*{s.}{Sobrenome: Di}
    \definition[个,位,种]{s.}{irmão mais novo | (entre amigos homens) eu | geralmente se refere a colegas do sexo masculino mais jovens na família ou entre parentes | forma humilde que os amigos usam para se referir uns aos outros, usada principalmente em correspondência}
  \antonymref{兄}{xiong1}
  \antonymref{长}{zhang3}
  \end{Phonetics}
\end{Entry}

\begin{Entry}{弟子}{7,3}{⼸,⼦}
  \begin{Phonetics}{弟子}{di4zi3}[][HSK 7-9]
    \definition{s.}{discípulo; pupilo; seguidor; nome antigo para estudante, aprendiz}
  \end{Phonetics}
\end{Entry}

\begin{Entry}{弟兄}{7,5}{⼸,⼉}
  \begin{Phonetics}{弟兄}{di4xiong5}
    \definition[位,伙,帮,群,个]{s.}{irmãos; irmão mais novo e irmão mais velho}
  \seealsoref{兄弟}{xiong1di5}
  \end{Phonetics}
\end{Entry}

\begin{Entry}{弟弟}{7,7}{⼸,⼸}
  \begin{Phonetics}{弟弟}{di4di5}[][HSK 1]
    \definition[个,位,名,些]{s.}{irmão mais novo | primo}
  \synonymref{兄弟}{xiong1di5}
  \antonymref{哥哥}{ge1ge5}
  \end{Phonetics}
\end{Entry}

\begin{Entry}{弟妹}{7,8}{⼸,⼥}
  \begin{Phonetics}{弟妹}{di4mei4}
    \definition{s.}{irmão e irmã mais novos | cunhada; esposa do irmão mais novo; usado principalmente na linguagem falada}
  \end{Phonetics}
\end{Entry}

%%%%%%%%%% 张 %%%%%%%%%%
\subsection*{张}\addcontentsline{loh}{figure}{张}

\begin{Entry}{张}{7}{⼸}
  \begin{Phonetics}{张}{zhang1}[][HSK 3]
    \definition*{s.}{Zhang, uma das vinte e oito constelações | Zhang, uma das mansões lunares | Sobrenome: Zhang}
    \definition{adj.}{indulgente; desenfreado; devasso; libertino}
    \definition{clas.}{usado para papel, couro, etc. | usado para cama, mesa, etc. | usado para o rosto, boca, etc. | usado para arco}
    \definition{v.}{abrir; espalhar; esticar | expor; exibir | (de uma loja) iniciar atividades comerciais; abrir | olhar; contemplar | expandir; estender; ampliar; exagerar | fixar (uma corda de arco); encordoar (um instrumento musical); puxar a corda do arco}
  \synonymref{绷}{beng1}
  \synonymref{紧}{jin3}
  \synonymref{开}{kai1}
  \antonymref{驰}{chi2}
  \antonymref{合}{he2}
  \antonymref{松}{song1}
  \end{Phonetics}
\end{Entry}

\begin{Entry}{张三}{7,3}{⼸,⼀}
  \begin{Phonetics}{张三}{zhang1san1}
    \definition*{s.}{Zhang San | Zé Ninguém | Nome para uma pessoa não especificada, 1 de 3}
  \seealsoref{李四}{li3si4}
  \seealsoref{王五}{wang2wu3}
  \end{Phonetics}
\end{Entry}

\begin{Entry}{张扬}{7,6}{⼸,⼿}
  \begin{Phonetics}{张扬}{zhang1yang2}[][HSK 7-9]
    \definition{v.}{divulgar; tornar público; dar ampla visibilidade; divulgar algo que é secreto ou que não precisa ser conhecido pelo público; espalhar rumores ou propaganda}
  \synonymref{高调}{gao1diao4}
  \synonymref{宣扬}{xuan1yang2}
  \antonymref{隐瞒}{yin3man2}
  \end{Phonetics}
\end{Entry}

\begin{Entry}{张灯结彩}{7,6,9,11}{⼸,⽕,⽷,⼺}
  \begin{Phonetics}{张灯结彩}{zhang1deng1-jie2cai3}[][HSK 7-9]
    \definition{expr.}{decorar com lanternas e serpentinas coloridas; enfeitar com lanternas e bandeirolas; decorado e adornado com lanternas; alegre com lanternas e decorações; festivo}
  \end{Phonetics}
\end{Entry}

\begin{Entry}{张狂}{7,7}{⼸,⽝}
  \begin{Phonetics}{张狂}{zhang1kuang2}
    \definition{adj.}{leviano e insolente; descarado | impetuoso | frenético}
  \synonymref{狂妄}{kuang2wang4}
  \synonymref{嚣张}{xiao1zhang1}
  \antonymref{低调}{di1diao4}
  \end{Phonetics}
\end{Entry}

\begin{Entry}{张南庄}{7,9,6}{⼸,⼗,⼴}
  \begin{Phonetics}{张南庄}{zhang1 nan2zhuang1}
    \definition*{s.}{Zhang Nanzhuang (anos de nascimento e morte desconhecidos), que viveu durante os períodos Qianlong e Jiaqing da Dinastia Qing, era conhecido como ``Passante'' (过路人); ele era bom em caligrafia no estilo de Ouyang Xun e imitava Fan Chengda e Lu You na escrita de poesia; ele gastava milhares de moedas de ouro todos os anos para colecionar livros raros, e sua coleção era a maior de Shanghai (上海) naquela época}
  \seealsoref{范成大}{fan4 cheng2da4}
  \seealsoref{过路人}{guo4lu4ren2}
  \seealsoref{陆游}{lu4 you2}
  \seealsoref{欧阳询}{ou1yang2 xun2}
  \seealsoref{上海}{shang4hai3}
  \end{Phonetics}
\end{Entry}

\begin{Entry}{张贴}{7,9}{⼸,⾙}
  \begin{Phonetics}{张贴}{zhang1tie1}[][HSK 7-9]
    \definition{v.}{postar; afixar; publicar}
  \end{Phonetics}
\end{Entry}

%%%%%%%%%% 弥 %%%%%%%%%%
\subsection*{弥}\addcontentsline{loh}{figure}{弥}

\begin{Entry}{弥}{8}{⼸}
  \begin{Phonetics}{弥}{mi2}
    \definition*{s.}{Sobrenome: Mi}
    \definition{adj.}{cheio; inteiro}
    \definition{adv.}{Literário: mais; ainda mais}
    \definition{v.}{transbordar; encher | cobrir; encher}
  \end{Phonetics}
\end{Entry}

\begin{Entry}{弥补}{8,7}{⼸,⾐}
  \begin{Phonetics}{弥补}{mi2bu3}[][HSK 7-9]
    \definition{v.}{remediar; compensar; reparar; restituir; consertar; preencher; inventar}
  \end{Phonetics}
\end{Entry}

\begin{Entry}{弥漫}{8,14}{⼸,⽔}
  \begin{Phonetics}{弥漫}{mi2man4}[][HSK 7-9]
    \definition{v.}{encher; transbordar; preencher o ar; espalhar"-se por toda parte}
  \end{Phonetics}
\end{Entry}

%%%%%%%%%% 弦 %%%%%%%%%%
\subsection*{弦}\addcontentsline{loh}{figure}{弦}

\begin{Entry}{弦}{8}{⼸}
  \begin{Phonetics}{弦}{xian2}[][HSK 7-9]
    \definition*{s.}{Sobrenome: Xian}
    \definition[根]{s.}{corda do arco; corda | Dialeto: mola (de um relógio, etc.) | corda de um instrumento musical | acorde | Matemática: hipotenusa}
  \end{Phonetics}
\end{Entry}

%%%%%%%%%% 弯 %%%%%%%%%%
\subsection*{弯}\addcontentsline{loh}{figure}{弯}

\begin{Entry}{弯}{9}{⼸}
  \begin{Phonetics}{弯}{wan1}[][HSK 4]
    \definition{adj.}{curvo; tortuoso; torto | para algo curvo, como a lua, etc. | dobrado; flexível}
    \definition[个,道]{s.}{curva; dobra; volta}
    \definition{v.}{dobrar; flexionar; curvar | Literário: desenhar}
  \synonymref{曲}{qu1}
  \antonymref{直}{zhi2}
  \end{Phonetics}
\end{Entry}

\begin{Entry}{弯曲}{9,6}{⼸,⽈}
  \begin{Phonetics}{弯曲}{wan1qu1}[][HSK 6]
    \definition{s.}{torto; curvo; sinuoso; tortuoso; não reto}
    \definition{v.}{dobrar; curvar; flexionar}
  \synonymref{波折}{bo1zhe2}
  \synonymref{曲折}{qu1zhe2}
  \antonymref{挺拔}{ting3ba2}
  \antonymref{挺立}{ting3li4}
  \end{Phonetics}
\end{Entry}

%%%%%%%%%% 弱 %%%%%%%%%%
\subsection*{弱}\addcontentsline{loh}{figure}{弱}

\begin{Entry}{弱}{10}{⼸}
  \begin{Phonetics}{弱}{ruo4}[][HSK 4]
    \definition{adj.}{fraco; debilitado | jovem | inferior; pior | colocado depois de uma fração ou decimal para indicar que é um pouco menor que esse número}
    \definition{v.}{perder (através da morte)}
  \synonymref{差}{cha4}
  \synonymref{故}{gu4}
  \synonymref{疲}{pi2}
  \synonymref{亡}{wang2}
  \synonymref{虚}{xu1}
  \synonymref{逊}{xun4}
  \synonymref{幼}{you4}
  \antonymref{刚}{gang1}
  \antonymref{健}{jian4}
  \antonymref{老}{lao3}
  \antonymref{强}{qiang2}
  \antonymref{生}{sheng1}
  \antonymref{胜}{sheng4}
  \antonymref{优}{you1}
  \antonymref{长}{zhang3}
  \antonymref{壮}{zhuang4}
  \end{Phonetics}
\end{Entry}

\begin{Entry}{弱势}{10,8}{⼸,⼒}
  \begin{Phonetics}{弱势}{ruo4shi4}[][HSK 7-9]
    \definition{s.}{uma tendência de queda; o ímpeto está diminuindo ou enfraquecendo | fraqueza; forças fracas}
  \end{Phonetics}
\end{Entry}

\begin{Entry}{弱点}{10,9}{⼸,⽕}
  \begin{Phonetics}{弱点}{ruo4dian3}[][HSK 7-9]
    \definition[个,种]{s.}{fraqueza; ponto fraco; áreas inadequadas; áreas frágeis}
  \end{Phonetics}
\end{Entry}

\begin{Entry}{弱弱}{10,10}{⼸,⼸}
  \begin{Phonetics}{弱弱}{ruo4 ruo4}
    \definition{adj.}{fraco}
  \end{Phonetics}
\end{Entry}

%%%%%%%%%% 弹 %%%%%%%%%%
\subsection*{弹}\addcontentsline{loh}{figure}{弹}

\begin{Entry}{弹}{11}{⼸}
  \begin{Phonetics}{弹}{dan4}
    \definition{s.}{bola; pelota; pequenas bolas disparadas com um estilingue | bomba; bala; explosivos que podem ser lançados ou arremessados, com poder destrutivo e letal}
  \end{Phonetics}
  \begin{Phonetics}{弹}{tan2}[][HSK 5]
    \definition{v.}{enviar; atirar (como com uma catapulta, etc.); usar a elasticidade de um objeto para lançar outro objeto | afofar; preparar fibras; usar um dispositivo elástico para amolecer as fibras | virar; sacudir | dedilhar; tocar (um instrumento musical de cordas) | acusar; atacar; criticar; relatar | saltar; quicar}
  \end{Phonetics}
\end{Entry}

\begin{Entry}{弹性}{11,8}{⼸,⼼}
  \begin{Phonetics}{弹性}{tan2xing4}[][HSK 7-9]
    \definition{adj.}{elástico; flexível; essa metáfora descreve a natureza das coisas, que são ajustáveis e adaptáveis de acordo com as necessidades reais}
    \definition{s.}{elasticidade; resiliência; a propriedade de um objeto se deformar sob a ação de uma força externa e retornar à sua forma original após a remoção dessa força}
  \synonymref{韧性}{ren4xing4}
  \end{Phonetics}
\end{Entry}

%%%%%%%%%% 强 %%%%%%%%%%
\subsection*{强}\addcontentsline{loh}{figure}{强}

\begin{Entry}{强}{12}{⼸}
  \begin{Phonetics}{强}{jiang4}
    \definition{adj.}{teimoso; inflexível}
  \end{Phonetics}
  \begin{Phonetics}{强}{qiang2}[][HSK 3]
    \definition*{s.}{Sobrenome: Qiang}
    \definition{adj.}{forte; poderoso | melhor; superior | mais; extra; adicional; um pouco mais que; usado após uma fração ou decimal para indicar que é um pouco maior que o número | resoluto; firme | violento | alto padrão}
    \definition{v.}{fortalecer; tornar forte; tornar poderoso}
  \seealsoref{棒}{bang4}
  \antonymref{弱}{ruo4}
  \end{Phonetics}
  \begin{Phonetics}{强}{qiang3}
    \definition{v.}{fazer um esforço; esforçar"-se}
  \end{Phonetics}
\end{Entry}

\begin{Entry}{强大}{12,3}{⼸,⼤}
  \begin{Phonetics}{强大}{qiang2da4}[][HSK 3]
    \definition{adj.}{forte; poderoso; potente; possante; descreve força forte e grande poder}
  \seealsoref{强盛}{qiang2sheng4}
  \synonymref{巨大}{ju4da4}
  \synonymref{庞大}{pang2da4}
  \synonymref{强劲}{qiang2jing4}
  \synonymref{强壮}{qiang2zhuang4}
  \synonymref{壮大}{zhuang4da4}
  \antonymref{崩溃}{beng1kui4}
  \antonymref{衰弱}{shuai1ruo4}
  \antonymref{微弱}{wei1ruo4}
  \antonymref{虚弱}{xu1ruo4}
  \end{Phonetics}
\end{Entry}

\begin{Entry}{强化}{12,4}{⼸,⼔}
  \begin{Phonetics}{强化}{qiang2hua4}[][HSK 6]
    \definition{v.}{intensificar; fortalecer; consolidar; tornar mais forte, melhorar sua habilidade e nível}
  \synonymref{巩固}{gong3gu4}
  \synonymref{加强}{jia1qiang2}
  \synonymref{深化}{shen1hua4}
  \synonymref{提高}{ti2//gao1}
  \synonymref{推进}{tui1jin4}
  \synonymref{增强}{zeng1qiang2}
  \antonymref{淡化}{dan4hua4}
  \antonymref{减轻}{jian3qing1}
  \antonymref{减弱}{jian3ruo4}
  \antonymref{削弱}{xue1ruo4}
  \end{Phonetics}
\end{Entry}

\begin{Entry}{强加}{12,5}{⼸,⼒}
  \begin{Phonetics}{强加}{qiang2jia1}[][HSK 7-9]
    \definition{v.}{forçar; impor; forçar alguém a aceitar uma determinada opinião ou prática}
  \end{Phonetics}
\end{Entry}

\begin{Entry}{强占}{12,5}{⼸,⼘}
  \begin{Phonetics}{强占}{qiang2zhan4}[][HSK 7-9]
    \definition{v.}{ocupar à força; apoderar"-se | ocupar à força; confiscar; anexar}
  \end{Phonetics}
\end{Entry}

\begin{Entry}{强壮}{12,6}{⼸,⼠}
  \begin{Phonetics}{强壮}{qiang2zhuang4}[][HSK 6]
    \definition{s.}{(corpo) forte, poderoso, robusto, resistente}
    \definition{v.}{fortalecer; construir}
  \synonymref{健康}{jian4kang1}
  \synonymref{健壮}{jian4zhuang4}
  \synonymref{结实}{jie1shi5}
  \synonymref{强大}{qiang2da4}
  \synonymref{雄壮}{xiong2zhuang4}
  \synonymref{壮大}{zhuang4da4}
  \synonymref{壮实}{zhuang4shi5}
  \antonymref{单薄}{dan1bo2}
  \antonymref{软弱}{ruan3ruo4}
  \antonymref{衰老}{shuai1lao3}
  \antonymref{衰弱}{shuai1ruo4}
  \antonymref{虚弱}{xu1ruo4}
  \end{Phonetics}
\end{Entry}

\begin{Entry}{强行}{12,6}{⼸,⾏}
  \begin{Phonetics}{强行}{qiang2xing2}[][HSK 7-9]
    \definition{adv.}{à força; com força}
  \end{Phonetics}
\end{Entry}

\begin{Entry}{强劲}{12,7}{⼸,⼒}
  \begin{Phonetics}{强劲}{qiang2jing4}[][HSK 7-9]
    \definition{adj.}{poderoso}
  \end{Phonetics}
\end{Entry}

\begin{Entry}{强制}{12,8}{⼸,⼑}
  \begin{Phonetics}{强制}{qiang2zhi4}[][HSK 7-9]
    \definition{v.}{forçar; compelir; coagir; utilizar o poder legal, político e econômico para forçar}
  \seealsoref{强迫}{qiang3po4}
  \synonymref{被迫}{bei4po4}
  \synonymref{逼迫}{bi1po4}
  \synonymref{迫使}{po4shi3}
  \synonymref{强行}{qiang2xing2}
  \synonymref{施压}{shi1ya1}
  \synonymref{压迫}{ya1po4}
  \synonymref{约束}{yue1shu4}
  \antonymref{自发}{zi4fa1}
  \antonymref{自由}{zi4you2}
  \antonymref{自愿}{zi4yuan4}
  \end{Phonetics}
\end{Entry}

\begin{Entry}{强势}{12,8}{⼸,⼒}
  \begin{Phonetics}{强势}{qiang2shi4}[][HSK 6]
    \definition*{adj.}{forte; poderoso; dominante}
    \definition{s.}{momento; ímpeto; grande impulso; forte impulso | força; influência dominante; forças poderosas}
  \synonymref{嚣张}{xiao1zhang1}
  \antonymref{劣势}{lie4shi4}
  \antonymref{弱势}{ruo4shi4}
  \end{Phonetics}
\end{Entry}

\begin{Entry}{强迫}{12,8}{⼸,⾡}
  \begin{Phonetics}{强迫}{qiang3po4}[][HSK 5]
    \definition{v.}{impelir; forçar; impor; compelir; aplicar pessão para obedecer}
  \seealsoref{逼迫}{bi1po4}
  \seealsoref{强制}{qiang2zhi4}
  \synonymref{勉强}{mian3qiang3}
  \synonymref{强行}{qiang2xing2}
  \synonymref{压迫}{ya1po4}
  \synonymref{压制}{ya1zhi4}
  \synonymref{抑制}{yi4zhi4}
  \antonymref{被迫}{bei4po4}
  \antonymref{随意}{sui2//yi4}
  \antonymref{自动}{zi4dong4}
  \antonymref{自发}{zi4fa1}
  \antonymref{自由}{zi4you2}
  \antonymref{自愿}{zi4yuan4}
  \antonymref{自主}{zi4zhu3}
  \end{Phonetics}
\end{Entry}

\begin{Entry}{强度}{12,9}{⼸,⼴}
  \begin{Phonetics}{强度}{qiang2du4}[][HSK 5]
    \definition[个,种]{s.}{intensidade; força | magnitude; rigor; avidez}
  \end{Phonetics}
\end{Entry}

\begin{Entry}{强项}{12,9}{⼸,⾴}
  \begin{Phonetics}{强项}{qiang2xiang4}[][HSK 7-9]
    \definition{adj.}{Literário: resoluto e inflexível; reto e inabalável}
    \definition{s.}{ponto forte (de um atleta ou de uma equipe) | jogo, evento ou assunto no qual alguém é forte | principal ponto forte | especialidade}
  \end{Phonetics}
\end{Entry}

\begin{Entry}{强烈}{12,10}{⼸,⽕}
  \begin{Phonetics}{强烈}{qiang2lie4}[][HSK 3]
    \definition{adj.}{muito forte; intenso; poderoso | violento; impetuoso; nível muito alto; atitude muito firme, sem espaço para mudanças | afiado; marcante; mostrado em contraste; muito claro}
  \seealsoref{激烈}{ji1lie4}
  \synonymref{剧烈}{ju4lie4}
  \synonymref{猛烈}{meng3lie4}
  \synonymref{强大}{qiang2da4}
  \synonymref{热烈}{re4lie4}
  \synonymref{鲜明}{xian1ming2}
  \antonymref{冷静}{leng3jing4}
  \antonymref{平淡}{ping2dan4}
  \antonymref{平和}{ping2he2}
  \antonymref{平静}{ping2jing4}
  \antonymref{平稳}{ping2wen3}
  \antonymref{轻微}{qing1wei1}
  \antonymref{柔和}{rou2he2}
  \antonymref{微弱}{wei1ruo4}
  \antonymref{温和}{wen1he2}
  \end{Phonetics}
\end{Entry}

\begin{Entry}{强调}{12,10}{⼸,⾔}
  \begin{Phonetics}{强调}{qiang2diao4}[][HSK 3]
    \definition{v.}{salientar; sublinhar; enfatizar; dar ênfase a; vincar}
  \synonymref{重申}{chong2shen1}
  \synonymref{加强}{jia1qiang2}
  \synonymref{突出}{tu1//chu1}
  \synonymref{着重}{zhuo2zhong4}
  \antonymref{淡化}{dan4hua4}
  \antonymref{忽略}{hu1lve4}
  \antonymref{轻视}{qing1shi4}
  \end{Phonetics}
\end{Entry}

\begin{Entry}{强盗}{12,11}{⼸,⽫}
  \begin{Phonetics}{强盗}{qiang2dao4}[][HSK 6]
    \definition[个,群,伙,帮]{s.}{ladrão; bandido; uma pessoa que usa violência para confiscar a propriedade de outros; também se refere a uma pessoa ou força que se envolve em comportamento semelhante}
  \synonymref{土匪}{tu3fei3}
  \antonymref{好人}{hao3ren2}
  \antonymref{君子}{jun1zi3}
  \end{Phonetics}
\end{Entry}

\begin{Entry}{强盛}{12,11}{⼸,⽫}
  \begin{Phonetics}{强盛}{qiang2sheng4}
    \definition{adj.}{(um país, movimento, etc.) poderoso e próspero; rico e poderoso}
  \seealsoref{强大}{qiang2da4}
  \synonymref{昌盛}{chang1sheng4}
  \synonymref{发达}{fa1da2}
  \synonymref{繁荣}{fan2rong2}
  \synonymref{蓬勃}{peng2bo2}
  \synonymref{强劲}{qiang2jing4}
  \synonymref{旺盛}{wang4sheng4}
  \synonymref{兴旺}{xing1wang4}
  \antonymref{没落}{mo4luo4}
  \antonymref{衰弱}{shuai1ruo4}
  \antonymref{衰退}{shuai1tui4}
  \antonymref{虚弱}{xu1ruo4}
  \end{Phonetics}
\end{Entry}

\begin{Entry}{强硬}{12,12}{⼸,⽯}
  \begin{Phonetics}{强硬}{qiang2ying4}[][HSK 7-9]
    \definition{adj.}{forte; resistente; inflexível; poderoso; não disposto a recuar}
  \end{Phonetics}
\end{Entry}

%%%%%%%%%% 彀 %%%%%%%%%%
\subsection*{彀}\addcontentsline{loh}{figure}{彀}

\begin{Entry}{彀}{13}{⼸}
  \begin{Phonetics}{彀}{gou4}
    \definition{adj.}{suficiente; adequado}
    \definition{v.}{puxar um arco ao máximo}
  \end{Phonetics}
\end{Entry}

%%%%% EOF %%%%%

